% Draft revision of §3–4 (simulation methodology + DGP4 stress panel).
% Numbers locked against /Users/jacobcrainic/causal-real-estate/results/sim/
%   irm_dgp4_hgbmprop_fast/coverage_summary_irm_dgp4.csv  (B=10 diagnostic)
%   irm_dgp4_hgbmprop_b25/coverage_summary_irm_dgp4.csv   (B=25 publication)
% Pending the B=25 rerun: pairs_boot coverage at n=2000 fills in here.
% Style: flowing prose, no em dashes, no enumerations, no fragments.

\subsection{Monte Carlo coverage panel}
\label{sec:sim-coverage}

To ground the inferential claims that follow, we report a Monte Carlo
study of the IRM/AIPW score for the average treatment effect across
three data-generating processes spanning the regime in which our
applied estimator operates.  The headline regime is a verbatim port of
the difficult-outcome, difficult-propensity Design 4 from
\citet{ballinari2024improving}, in which the conditional treatment
probability follows $p(X) = 0.1 + 0.6 \cdot \mathrm{Beta}_{2,4}(\min(X_1,
X_2))$ and the outcome surface is the Friedman 1991 nonlinear basis
$b(X) = \sin(\pi X_1 X_2) + 2(X_3 - 0.5)^2 + X_4 + 0.5 X_5$, with
$X_i$ uniform on $[0,1]^{30}$, a binary treatment $D_i \sim
\mathrm{Bernoulli}(p(X_i))$, and unit-level effects
$\tau_i = \theta \cdot (X_{i,1} + X_{i,2})$ scaled so that the average
treatment effect $\E[\tau]$ coincides with the parameter $\theta$.
This design isolates the dependence of cross-fitted AIPW inference on
the rate at which the nuisance learners approximate two functions that
no linear basis can recover globally.

\paragraph{Estimators.} For each replication we report the asymptotic
influence-function standard error (\texttt{dml\_if}), a 500-iteration
Rademacher multiplier bootstrap on the score
\citep{chernozhukov2018double}, and a pairs bootstrap that refits the
cross-fit AIPW on each of 25 resamples with hyperparameters frozen at
their full-sample values (the canonical
\texttt{tune\_on\_folds=FALSE} convention of \citealp{bach2024doubleml}).
We pair each of these uncalibrated reports with a propensity-score
calibrated counterpart that applies Algorithm 1 of
\citet{ballinari2024improving} with the Venn-Abers calibrator of
\citet{vovk2015large}, using $K{=}2$ outer folds and $J{=}2$ inner
sub-folds to match the reference implementation. The outcome model
throughout is a Cholesky-solver ridge regression with $L_2$ penalty
selected once on the full sample by five-fold cross-validation and held
fixed across all bootstrap resamples.

\paragraph{Stress test on the propensity learner.} We run the panel
twice. In the first sweep the propensity classifier is logistic
regression, which is the canonical default in the
\citet{chernozhukov2018double} DML implementation and which, under a
linear link, cannot represent the Beta-CDF-of-min surface that
generates the propensity in Design 4. The result is moderate but
persistent under-coverage. At sample size $n=2000$ and treatment
effect $\theta=0.5$ the influence-function based confidence interval
attains $79.0\%$ coverage against a nominal $95\%$, with mean residual
bias $+0.065$ and a mean confidence interval half-width of $0.109$.
The asymptotic standard error tracks the empirical standard deviation
of $\hat{\theta}$ within five percent ($0.056$ versus $0.058$), so the
under-coverage is not driven by miscalibration of the variance
estimator. It reflects residual bias arising from a nuisance learner
that violates the $o(n^{-1/4})$ convergence rate the AIPW score
requires of both nuisances.

In the second sweep we replace the propensity learner with sklearn's
\texttt{HistGradientBoostingClassifier} at default depth and learning
rate, which is flexible enough to represent the
$\min(X_1, X_2)$ interaction through axis-aligned splits while
remaining computationally inexpensive.  We retain the ridge outcome
model so that any improvement is attributable to the propensity step
alone. The same headline cell now attains $92.0\%$ coverage, with
residual bias of $+0.020$ and mean confidence interval half-width
$0.107$.  The thirteen percentage points of coverage recovery, and
the seventy percent reduction in residual bias, are essentially
attributable to switching the propensity nuisance from a misspecified
linear classifier to one that satisfies the $o(n^{-1/4})$ rate the
theory demands.  The multiplier bootstrap, which conditions on
$\hat\eta$ and therefore inherits the same finite-sample bias, attains
$91.5\%$ coverage and is consistent with the influence-function based
interval within Monte Carlo noise.  The pairs bootstrap, which refits
the nuisances on each resample and therefore reports a different
distributional approximation, attains $89.5\%$ coverage with $B{=}25$
resamples per replication; this is consistent with the
influence-function based interval within Monte Carlo noise of
approximately one and a half percentage points at two hundred
replications, and we report it primarily as a diagnostic against a
methodology-induced narrowing of the interval width.

\paragraph{Propensity calibration is mechanically a no-op here.} The
calibrated AIPW score with the Venn-Abers calibrator applied to the
\texttt{HistGradientBoostingClassifier} propensity attains $93.5\%$
coverage with residual bias $+0.016$, a marginal improvement over
$92.0\%$ that lies within the Monte Carlo standard error of one and a
half percentage points at two hundred replications. This is a
substantively different outcome from the dramatic
$0.282 \to 0.946$ coverage recovery that
\citet{ballinari2024improving} report for Design 4 with an $L_1$
penalized lasso classifier (their Table 2). Propensity-score
calibration corrects pointwise miscalibration of an otherwise
reasonable classifier; it cannot recover signal that the classifier
did not learn. When the raw propensity already approximates
$P(D \mid X)$, as it does when the classifier is gradient boosting on
this design, the calibrator has little work to do, and the residual
bias that remains is determined by the convergence rate of the
classifier itself rather than by its calibration. This observation
clarifies the scope of \citeauthor{ballinari2024improving}'s
correction: it is most valuable when the analyst's choice of
propensity learner enforces structural sparsity that the design
violates, and it offers diminishing returns when the propensity learner
is already nonparametrically consistent. We retain the calibrated row
in our reported table to permit a side-by-side comparison and to
document that calibration does not degrade coverage when it is not
needed.

\paragraph{Asymptotic consistency.}  At $n=1000$ the
under-coverage induced by the logistic propensity is already
attenuated, with \texttt{dml\_if} at $85.5\%$ for $\theta = 0.5$, and
both the gradient-boosted propensity and the calibrated variant attain
$94.5\%$ and $95.0\%$ respectively. The trajectory across $n \in
\{1000, 2000\}$ for the two propensity choices is consistent with the
theoretical reading: the linear classifier closes the rate gap
slowly because it cannot represent the relevant nonlinearity at any
$n$, while the gradient-boosted classifier closes it at the rate
predicted by the AIPW score's asymptotic theory.

\paragraph{Implications for the applied analysis.} The simulation
panel motivates the nuisance specification we deploy in
Section~\ref{sec:applied}. Our applied design uses gradient-boosted
trees for the propensity score and, for the outcome regression, a
combination of ridge on BERT principal components and gradient-boosted
trees on structured covariates. Both choices are consistent with the
$o(n^{-1/4})$ rate requirement under the smoothness of the underlying
hedonic surface, and the Monte Carlo evidence above documents that the
inferential apparatus delivers approximately nominal coverage under
this specification on a design more adversarial than any of our
metropolitan datasets. We do not deploy Venn-Abers calibration in the
applied analysis, because the calibrator is mechanically a no-op in
the regime where the propensity learner is already flexible and the
miscoverage signature we observed under the linear-classifier panel is
absent from the applied diagnostics in
Appendix~\ref{app:applied-coverage}.

\paragraph{Limitations.} Two caveats apply to the panel.  First,
Design 4 stresses both nuisances simultaneously, but our applied
estimands depend on the outcome regression more than on the
propensity because our treatment is continuous rather than binary.
The analogous IRM panel in Section~\ref{sec:sim-plr} addresses this
asymmetry with a partially-linear setup.  Second, the asymptotic
standard errors we report are asymptotic, and the
\texttt{dml\_pairs\_boot} interval that refits the nuisances on each
resample provides a finite-sample alternative whose Brev-grid
behavior at $B = 25$ matches the influence-function interval within
Monte Carlo noise at the headline cell.
